% =====================================================================
% VERSION INSTITUT A 150 PAGES : protocole d'elicitation prepare et relectures de praticiens.
% Reprise de supplement/S4 (sections « methode de Cooke » a « cout de l'absence »), reecrite en
% phrases courtes. Les nombres, etiquettes et references sont ceux de l'original. Les destinataires
% de l'appel et les repondants ne sont pas nommes. Lue par main_institut_150.tex seul.
% =====================================================================

\chapter{Protocole d'élicitation et relectures de praticiens}
\label{ch:elicitation}

Le protocole d'élicitation de la matrice de contagion a été préparé, diffusé, puis laissé sans panel.
Trois relectures qualitatives ont été reçues à la place.

\section{Méthode de Cooke}
% HARNAIS-HORS-SCRIPT: quantiles demandes a l'expert (5, 50 et 95 %) : propriete de conception du questionnaire, non une grandeur calculee.
% HARNAIS-HORS-SCRIPT: quantiles demandes a l'expert (5, 50 et 95 %) : propriete de conception du questionnaire, non une grandeur calculee.

Le modèle classique de Cooke \citep{Cooke1991} ne traite pas les experts comme interchangeables. Il les note sur des
questions dont l'analyste connaît la réponse, puis pondère leurs avis par cette note. La méthode se
déroule en quatre temps.
\begin{enumerate}[itemsep=3pt,topsep=3pt]
\item Chaque expert donne trois quantiles, à $5$, $50$ et $95\,\%$, plutôt qu'un point. Son incertitude
      déclarée devient mesurable.
\item Sur les questions-graines, dont la réponse est connue de l'analyste et masquée à l'expert, on
      compare la fréquence empirique des réalisations dans les intervalles déclarés à la fréquence
      théorique. L'écart se lit comme une statistique d'information de Kullback-Leibler, et son niveau
      de signification est le score de calibration.
\item Un expert bien calibré qui déclare des intervalles très larges n'apporte rien. Le score
      d'information mesure donc la concentration de ses distributions par rapport à une loi de référence.
\item Le poids d'un expert est le produit des deux scores, et il est nul sous un seuil de calibration.
      L'avis agrégé est la combinaison linéaire pondérée des distributions individuelles.
\end{enumerate}
Le quatrième point fait la force et la fragilité de la méthode. Un expert mal calibré est écarté, ce qui
protège des avis péremptoires, mais la taille effective du panel diminue d'autant.

\section{Protocole construit pour la matrice de contagion}

Le questionnaire comporte deux parties. La partie d'étalonnage pose des questions à réponse connue, tirées
de sources publiques vérifiables, sur des grandeurs du même domaine que les cibles. La partie cible
interroge un à un les liens dirigés entre piliers, sous la forme d'une probabilité conditionnelle : sachant
que le pilier source est défaillant, quelle est la probabilité que le pilier cible le devienne ?

Le dispositif complet demande seize questions par expert, dont dix questions-graines. Ce nombre est un choix de l'instrument, pas un minimum statistique universel. La stabilité des poids resterait à examiner sur les réponses recueillies. Le protocole, le questionnaire et le gabarit de réponses
sont versionnés dans le dépôt, et deux scripts implémentent l'agrégation : l'un sur des réponses fictives,
pour valider la chaîne de traitement, l'autre prêt à recevoir des réponses réelles.
% SOURCES-SCRIPTS: 26 47 63

Une graine a été retirée de la première version. Elle portait sur l'indice de queue de sévérité
$\xi$, de valeur vraie annoncée $0{,}90$. Or $0{,}90$ est la valeur posée du scénario de pire cas, la
calibration figée donnant $\hat\xi = 0{,}5954$ ; dans la méthode de Cooke, une graine fausse inverse
les poids en pénalisant le répondant qui vise juste. De plus, $\xi$ est une estimation dont
l'intervalle $[0{,}30\,;0{,}83]$ a une largeur proche de la valeur elle-même, alors qu'une graine doit
être une grandeur réalisée. Les dix graines de l'instrument révisé sont toutes reproductibles par une
sortie de script versionnée.

\section{Documents produits et diffusés}
% HARNAIS-HORS-SECTION: cette section rapporte une demarche et son issue, non une grandeur
% calculee. Ses comptes sont ecrits en lettres, aucun script ne les imprime et aucun ne le
% devrait : ce qui les atteste est le materiel verse au depot, pas une sortie numerique.

Quatre pièces sont versionnées au dépôt : le protocole, qui fixe la méthode, les seuils et la règle
d'agrégation ; le questionnaire, dans sa troisième version ; un formulaire en ligne reprenant ses
questions ; le gabarit de réponses, qui alimente la chaîne d'agrégation. L'appel a été adressé à une
commission de l'Institut des Actuaires et à des praticiens du cabinet, qui ne sont pas nommés. Il est
resté sans réponse, et aucun panel n'a pu être constitué. Ce fait de démarche n'est pas présenté comme
un résultat : il dit la difficulté d'obtenir un jugement quantitatif structuré sur un objet aussi
récent, et rien de la disposition des personnes sollicitées.

\section{Portée de l'absence de panel}

L'absence de panel limite la confrontation externe du barème. Une à deux réponses auraient aussi
laissé un panel fragile : un seul expert pourrait franchir le seuil, ou les poids dépendre fortement
des dix graines. Cela ne démontre pas qu'un petit panel serait nécessairement sans information.

Une élicitation biaisée pourrait aussi être trompeuse. Dans l'expérience, un panel biaisé déplace le SCR de $777$~M\euro{}
sans qu'aucun signal n'avertisse le lecteur (chapitre~\ref{chap:partielle}), alors que l'approche par
bornes reste invariante si seule la direction change et si la partie symétrique d'expert reste fixée. Cette invariance ne protège pas d'un biais sur la force des liens : sans panel, celle-ci demeure posée.
% SOURCES-SCRIPTS: 31

Ce travail s'appuie pourtant ailleurs sur un jugement d'expert : le codage, par un seul codeur, du
corpus de rapports post-incident. Trois différences le justifient. Le codage est auditable, car le
protocole est écrit et les récits sont publics. Son biais principal est testé contre une loi nulle
exacte (\S\ref{sec:biais-narration}). Son produit est un ordre partiel, que la suite borne, et non des
coefficients qui entreraient directement dans le capital. Le codage à juge unique reste une limite
déclarée, bornée au \S\ref{sec:biais-narration} ; un second codage en aveugle la réduirait.

\section{Dispositifs de substitution}

Trois dispositifs complètent le jugement d'expert, sans remplacer la partie symétrique posée.
\begin{itemize}[itemsep=3pt,topsep=3pt]
\item Les rapports post-incident officiels sont codés selon un protocole écrit, et le sens des
      transitions est testé contre une loi nulle exacte de permutation ($z = +5{,}12$ sur le corpus
      étendu).
\item Une échelle de repli suppose la direction de rang faible. Son premier cran capte $74{,}1\,\%$
      de la variance de la direction posée et resserre le capital dans $[7\,761\,;\,8\,316]$~M\euro,
      soit une largeur de $555$~M\euro{} contre $1\,839$ sous ignorance totale. C'est une hypothèse
      structurelle, déclarée comme telle.
\item Les niveaux d'exigence des textes et les publications des autorités fournissent des repères
      externes sur les états de conformité.
\end{itemize}
% SOURCES-SCRIPTS: 30 34 59

\section{Relectures de praticiens}
\label{sec:relecture-praticien}

La direction de $W$ n'étant pas identifiable sur les données disponibles
(chapitre~\ref{ch:identifiabilite}), un seul contrôle externe reste possible : le jugement de
praticiens qui ont vu des défaillances réelles. Un praticien peut corroborer ou contredire une
structure, non la calibrer. L'instrument soumet donc sept affirmations qualitatives du modèle, dont
deux sur l'ordre de P1 vers P2, en demandant de les contredire : un désaccord argumenté désigne
l'endroit où le modèle s'écarte du terrain, un accord ne discrimine rien.

\emph{Première réponse.} Une consultante de Nexialog, en gestion des risques et du contrôle depuis
une dizaine d'années, ne contredit aucune des sept phrases (deux accords pleins, cinq nuancés). Elle
juge les trois indicateurs pilotables cohérents et produisibles, et décrit l'usage établi d'un
tableau de bord des incidents par prestataire pour évaluer, au titre de DORA, la fréquence de
défaillance système et le taux d'indisponibilité sur douze mois. Sa nuance principale distingue
l'occurrence et la maîtrise : une gouvernance solide ne change pas le fait qu'un prestataire critique
tombe, mais la capacité à en limiter les conséquences. Dans le modèle, cela déplacerait l'arc de la
gouvernance vers les tiers de la matrice de propagation vers le canal de détection. Cette réserve est
déclarée sans être traitée, le gel interdisant la recalibration. Sur la gestion des incidents comme
puits, elle avance un argument de précédence (les tests précèdent la mise en service), que le corpus
ne code pas. Cette réponse a conduit à ajouter trois questions à choix forcé, dont l'une offre la
modalité \emph{je ne peux pas trancher}, qui corroborerait la frontière d'identifiabilité.

\emph{Deuxième réponse.} Un consultant de Nexialog, expert de la résilience dans les services
financiers depuis sept ans, donne quatre accords nuancés et deux désaccords argumentés. Il note que
plusieurs domaines peuvent céder ensemble par une cause partagée (manque de ressources,
transformation mal pilotée, architecture vieillissante) : c'est la frontière d'identifiabilité
formulée par un praticien. Il conteste que la gouvernance soit une source pure et la gestion des
incidents un puits. La plupart de ses exemples décrivent une rétroaction corrective, qui relève de
l'évolution des états de conformité (chapitre~\ref{chap:multietats}). Un seul décrit une dégradation
en retour, de la gestion des incidents vers le partage d'information, sens qu'aucun rapport du corpus
ne code. Ce désaccord ne déplace pas le capital publié, l'ensemble admissible du
chapitre~\ref{chap:partielle} ne fixant aucune direction ; il affaiblit la structure suggérée par le
corpus, déjà déclarée corroborée et non calibrée. Il juge le registre réaliste sous condition de règles
de qualification précises, et signale un manque : la gestion des changements et des vulnérabilités,
déclencheurs transversaux sans place propre dans le modèle.

\emph{Troisième réponse.} Un consultant en données, qui exerce depuis un an, conteste les deux mêmes
phrases et juge le registre réaliste mais risqué pour la confidentialité. Il n'apporte pas d'argument
nouveau et a demandé à n'être cité que par sa fonction.

Trois réponses ne font pas un panel. Cette relecture n'entre dans aucune table de robustesse et ne
modifie aucun niveau publié.

\section{Coût de l'absence et réversibilité}

L'élicitation devait départager P1 et P4 en tête du classement des piliers sources. Sans elle, les
deux restent presque à égalité sur $220$ sommets admissibles tirés au hasard, P1 arrivant premier dans
$43{,}6\,\%$ des tirages et P4 dans $37{,}7\,\%$, pour un écart de regret maximal observé de $1{,}5\,\%$
(\S\ref{sec:res-robustesse}). P1 et P4 devancent nettement les trois autres piliers, et l'analyse ne
tranche pas entre eux par un avis.
% SOURCES-SCRIPTS: 30 31

Le protocole, le questionnaire, le gabarit et la chaîne d'agrégation sont prêts et testés ; il manque
un panel d'une dizaine de répondants. La reprise demanderait recrutement et contrôle des réponses. Elle pourrait départager P1 et P4 et resserrer la bande, sans gain garanti à l'avance.
